We study how reliably automated machine learning (AutoML) selects anomaly detection models for Internet of Things (IoT) systems, and we show that the evaluation protocol influences the reported performance more than the choice of framework. We compare the widely used frameworks Auto-sklearn and H2O AutoML on public sensor and network-traffic datasets under identical data, search budgets, thread counts, and random seeds, with default models as a reference. We propose AutoML-IoT, an evaluation pipeline that audits the data, constructs leakage-aware protocols, runs controlled searches, measures detection quality and deployment cost, and explains the selected models. We provide empirical evidence that holding out entire recording blocks lowers the score of every method, whereas the difference between the frameworks stays within the variation across random seeds. We also show that duplicated and label-ambiguous records, which cover 40.5\% of the rows of a traffic dataset, bound the attainable accuracy at 0.72. We then identify the effects of class balancing, data cleaning, and split strategy on the ranking of the frameworks. Both frameworks perform close to a default random forest, but they differ substantially in model size and prediction latency. Permutation importance indicates that both frameworks rely on the same few traffic statistics. We release code and per-run results to support replication.
